\begin{tablalafrox}{Diccionario de la cuenta 4.3. Fuente: (Bases Administrativas, arts. 17.2–17.6 y Formulario T-14, p. 63); (Bases Técnicas Transversales, secciones 17.2–17.6, pp. 32–34).}{tab:T-14-31}
  {L{1.2cm} Y Y Y C{0.9cm} L{2.3cm}}
  {\cab{Código} & \cab{Paquete de trabajo} & \cab{Entregable} & \cab{Criterio de aceptación} & \cab{Resp.} & \cab{Período}}
  4.3.1 & Marcha blanca de la Etapa 2 en convivencia con la Etapa 1 y fuente única de datos (H11) & Informes diarios de indicadores y de conciliación; informe mensual de restauración de respaldos; acta de cierre de la marcha blanca. & Se cumplen las seis condiciones del Art. 17.3, no hay doble digitación y cada mes se restaura una muestra representativa sin errores con su tiempo medido (RT-07.12). La Contraparte Técnica firma el acta de cierre tras revisar la evidencia del mes 20, lo que habilita el paso a producción del mes 21. & IMP & Meses 19 y 20. Empieza con el H11. \\
  4.3.2 & Estabilización de la Etapa 2 con Comercial, Finanzas, cadenas y transportistas & Informe de estabilización. & No quedan defectos críticos ni altos abiertos al cierre del período. & IMP & Desde el paso a producción de la Etapa 2 hasta el cierre del mes 22, con cuatro semanas reforzadas posteriores a la firma del H12; con inicio contractual supuesto en febrero de 2027, del 5 de octubre al 30 de noviembre de 2028, con refuerzo del 26 de octubre al 22 de noviembre. \\
  4.3.3 & Expediente de H12 & Expediente con evidencia de 4.3.1 y registro de producción. & Firma previa al corte y registro de producción efectiva en el mes 21 con garantía constituida; secuencia conforme al T-15, sección 5.5. & JP & Preparación durante la marcha blanca; aceptación en el mes 21. \\
  4.3.4 & Garantía de correcto funcionamiento constituida, condición de la aceptación final (H12) & Garantía constituida. & Está entregada antes de la entrega del expediente del H12. & JP & Mes 20, en paralelo con la marcha blanca de la Etapa 2; es un trámite financiero y no interviene sistemas. \\
\end{tablalafrox}
